\section{Results}
\label{sec:results}

\subsection{Placement on the block}
\label{subsec:block}

\begin{figure*}[t]
  \centering
  \includegraphics[width=0.88\textwidth]{f2b-crossover-vs-d.png}
  \caption{Split-update time divided by replicated-update time on the block benchmark.
    Below 1, splitting is faster; above 1, replication is faster. Panels separate hardware
    and matrix width. Marker shape identifies population size. Bars show the range from
    fastest split / slowest replicated repeat to slowest split / fastest replicated
    repeat, across five repeats per placement; these are not confidence intervals.
    The three main variants measured on both platforms are shown; the complete grids,
    including pairing variants, are in Appendix~\ref{sec:grids}.}
  \label{fig:f2}
\end{figure*}

At eight devices in one host, splitting the update is faster for every full-rank
configuration. For the dense and seed variants drawn in Figure~\ref{fig:f2}, the speedup
is 1.45--5.62$\times$. Dense alone gains 1.8--2.4$\times$ on the A100 and
5.0--5.6$\times$ on the v5e, so this result does not depend on the slower seed-regeneration
evaluation path. Replication is faster for rank-1 perturbations at width 2048
(Figure~\ref{fig:f2}). For those low-rank configurations, splitting takes 4--15\% longer
on the A100 and 17--46\% longer on the v5e. These low-rank ranges include both mirrored (drawn) and
unpaired sampling; Appendix~\ref{sec:grids} gives all variants and absolute times.

At width 512 the rank-1 placements are close. We treat a comparison as unresolved when
its repeat range contains 1, using the range defined in the figure caption. At eight
devices this occurs in six of the eight low-rank configurations at width 512. Every
full-rank configuration and every low-rank configuration at width 2048 has a repeat
range entirely on one side of 1. These repeats measure variation within a compiled run,
not variation across hardware sessions.

% generated by experiments/phase2/paper_evidence.py; do not edit
\begin{table*}[t]
\centering
\small
\caption{Worked example of Eq.~\ref{eq:crossover}, at width 2048, population 256 and eight devices. All times are milliseconds. $C_A$ and $C_B$ are separately measured reconstruction costs; $T_{\mathrm{add}}$ is communication added to the split reconstruction probe. The prediction is $C_A-C_B+T_{\mathrm{ag}}-T_{\mathrm{add}}$, including the small weight gather. Positive differences favor splitting. $t_A$ gives the generation time against which a penalty is divided in Figure~\ref{fig:f2}.}
\label{tab:worked}
\begin{tabular}{llrrrrrr}
\toprule
platform & variant & $C_A$ & $C_B$ & $T_{\mathrm{add}}$ & predicted $t_A-t_B$ & measured $t_A-t_B$ & $t_A$ \\
\midrule
A100 & seed & 61.05 & 7.78 & 1.25 & 52.03 & 54.30 & 113.26 \\
A100 & rank 1 & 0.77 & 0.35 & 1.62 & -1.19 & -1.40 & 19.88 \\
v5e & seed & 233.29 & 29.51 & 2.03 & 201.75 & 213.30 & 467.58 \\
v5e & rank 1 & 0.55 & 0.33 & 2.36 & -2.14 & -2.04 & 9.73 \\
\bottomrule
\end{tabular}
\end{table*}


Table~\ref{tab:worked} evaluates Eq.~\ref{eq:crossover} at one configuration on each
platform. Full-rank reconstruction costs 61--233\,ms, so splitting it easily pays for
1--2\,ms of communication. Rank-1 reconstruction costs less than 1\,ms; the communication
is then more expensive than the work saved. On the A100, the model gives a 1.19\,ms
penalty for splitting rank 1, against 1.40\,ms measured. On the v5e the measured penalty
is 2.04\,ms on a shorter replicated generation: 9.73\,ms against 19.88\,ms on the A100.
That makes the relative penalties 21\% and 7\% respectively.

The ratio plots and the timing equation describe the same difference:
$t_B/t_A=1-(t_A-t_B)/t_A$. Thus a fixed communication penalty can look larger as a ratio
when evaluation becomes faster. Table~\ref{tab:tb8} in the appendix reports the isolated communication
cost; Table~\ref{tab:worked} instead uses the cost added to a reconstruction probe,
which also includes any scheduling changes that operation causes.

\begin{figure*}[t]
  \centering
  \includegraphics[width=\textwidth]{f9-e17-crossover.png}
  \caption{Placement on Qwen2.5-0.5B, TPU v5e: split / replicated time for a complete
    generation, including next-token loss evaluation. Splitting is faster for full-rank
    noise; replication is faster at every retained low-rank configuration. The strips
    below the plots distinguish fits, memory failures (OOM), and excluded comparisons.
    Only records stamped with a clean worktree are used; gaps are not connected.
    All one-device records are excluded. Absolute times are in Table~\ref{tab:tb5}.}
  \label{fig:f9}
\end{figure*}

Using separately measured reconstruction and added communication costs, the model
agrees with all 30 configurations whose repeat envelopes do not overlap. It agrees in
33 of 36 overall; its three misses are among the six near-ties. This is a check of whether
separately timed components account for complete generations, not a fit to their timing
differences. It was developed after the sweeps. Assuming ideal $1/D$ reconstruction and
using only the isolated collective benchmark misses two resolved A100 configurations.
Appendix~\ref{sec:timemodel} gives the comparison and the remaining errors in magnitude.

\subsection{Placement on a real model}



The same preference holds on Qwen2.5-0.5B on the TPU (Figure~\ref{fig:f9}).
We retain only records stamped with a clean worktree, leaving eight full-rank and twelve
low-rank timing comparisons. All eight-device records pass this filter.
Splitting helps full rank, with a larger advantage at eight than four devices.
Replication helps low rank; the rank-16, $N=128$ comparison is nearly tied, with a ratio
of 1.009 and a repeat envelope starting at 1.0006. Appendix~\ref{sec:memory} documents
the excluded records rather than treating them as missing because of memory limits.

The update contains 494,032,768 float32 values, including the tied embedding once:
1.84\,GiB. Scaling the isolated v5e communication time at 96\,MiB to this payload gives
39.4\,ms. This is an extrapolation of the measured communication rate, with no parameter
fit to Qwen's timing. The extra split cost for ranks 1 and 4 is 34--48\,ms in the
retained four- and eight-device comparisons. The payload estimate therefore explains
its scale, though it does not isolate the reconstruction contribution or test the
collective at this larger payload directly.

At eight devices, the rank-1 penalty stays roughly fixed as $N$ grows, while evaluation
time grows and the ratio approaches 1. Conversely, adding devices shortens evaluation,
so a nearly fixed penalty can produce a rising ratio. Rank 16 also reduces the absolute
penalty, to 2.4\,ms at $N=128$. The rank trend is not monotone: at $N=64$, rank 4 has a
ratio of 1.40, above rank 1's 1.31. The signs and a smaller rank-16 margin were expected
before this grid, after an earlier two-variant study; the population trend was observed
during the run.

At eight devices, all low-rank variants run out of memory at $N=240$, while full rank
fits. Low-rank evaluation batches candidates together; the full-rank variant regenerates
and evaluates them one at a time. This trades activation and weight memory against
throughput. Reduced-model compilation probes support the activation-memory explanation
(Appendix~\ref{sec:memory}); they do not establish an identical memory boundary on Qwen
with matched batching. Each placement comparison keeps batching fixed, whereas a timing
comparison across perturbation types includes this evaluation difference.

\subsection{Across a slow connection between hosts}
\label{subsec:boundary}

% generated by experiments/phase2/multihost/tb7_e18.py --latex; do not edit
\begin{table*}[t]
\centering
\small
\caption{The placement across a host boundary: $t_B - t_A$ in ms, negative where all-reduce (B) wins. Two 8$\times$A100 nodes joined by TCP sockets (InfiniBand carried no NCCL traffic); \texttt{1x8} is one node at $D{=}8$, \texttt{2x4} the same eight devices split across the boundary, \texttt{2x8} sixteen devices across it. The all-reduce runs at 64\,GiB/s on NVLink and 0.82 (\texttt{2x4}) and 0.58\,GiB/s (\texttt{2x8}) across the boundary. \emph{prereg.}: the prediction committed before the \texttt{2x8} configurations ran, whose bandwidth was $D$ times too high (the calibration all-reduced $1/D$ of its label). \emph{model} and \emph{corrected}: Eq.~\eqref{eq:crossover} on the same inputs with the bandwidth at the payload actually moved, computed after the measurements. The \texttt{1x8} column is re-measured on the cluster's first node; the model starts from the single-node sweep.}
\label{tab:tb7}
\begin{tabular}{llrrrrrrr}
\toprule
 & & & & \multicolumn{2}{c}{\texttt{2x4}} & \multicolumn{3}{c}{\texttt{2x8}} \\
\cmidrule(lr){5-6}\cmidrule(lr){7-9}
arm & $d$ & $N$ & \texttt{1x8} & measured & model & measured & prereg. & corrected \\
\midrule
seed & 2048 & 128 & $-21.4$ & $+86.8$ & $+91.0$ & $+74.1$ & $-13.9$ & $+137.2$ \\
seed & 2048 & 256 & $-50.8$ & $+57.2$ & $+58.7$ & $+56.7$ & $-48.2$ & $+103.0$ \\
seed & 512 & 1024 & $-20.1$ & $-13.0$ & $-18.9$ & $-24.8$ & $-27.0$ & $-17.6$ \\
rank 1 & 2048 & 128 & $+1.6$ & $+108.9$ & $+114.6$ & $+159.1$ & $+11.7$ & $+162.7$ \\
rank 1 & 2048 & 256 & $+1.4$ & $+111.4$ & $+114.4$ & $+130.9$ & $+11.4$ & $+162.5$ \\
rank 1 & 512 & 1024 & $0.0$ & $+6.9$ & $+7.1$ & $+7.8$ & $+0.8$ & $+10.2$ \\
\bottomrule
\end{tabular}
\end{table*}


Communication can reverse the full-rank preference. We repeated six configurations on
two A100 nodes. NCCL could not use the provider's InfiniBand, so cross-host traffic
ran over TCP sockets. The calibrated effective cross-host all-reduce bandwidth was
0.6--0.8\,GiB/s.

Keeping eight devices but splitting them between hosts adds 107--110\,ms to the
split placement's relative cost at width 2048 (Table~\ref{tab:tb7}). This exceeds the
21--51\,ms it saved on one host, so replication becomes faster for regenerated full-rank
noise. At width 512, the update is only 6\,MiB and splitting remains faster for that
variant. Replication's advantage for rank 1 grows across the slow connection.

The rule also gives a useful threshold for another connection. Keeping eight devices
and neglecting latency, the 96\,MiB all-reduce stops paying below about 4.1\,GiB/s at
$N=128$ and 1.8\,GiB/s at $N=256$. These are effective all-reduce rates, not network
line rates. They use the within-host savings plus the within-host communication being
replaced; Appendix~\ref{sec:boundary-audit} gives the calculation and repeat bounds.

Our prediction fixed before the sixteen-device run identified four of six faster
placements. Its bandwidth calibration used an incorrectly labelled message size.
Correcting that error after the runs gives all six preferences, but still overestimates
the split penalty by up to 63\,ms. The original prediction and the
retrospective correction remain side by side in Table~\ref{tab:tb7-audit}.
This experiment establishes the effect of a slow connection, not the performance of a
high-bandwidth multi-host cluster.

\subsection{What low-rank perturbations save}

\begin{figure*}[t]
  \centering
  \includegraphics[width=0.9\textwidth]{f4-cost-common-bfloat16.png}
  \caption{Generation time relative to stored full-rank noise (dense), on one device
    with bfloat16 storage and default matrix-multiplication precision. Every row uses
    the same model width and population on both platforms, and all compared variants fit.
    Values below 1 are faster than dense. Low-rank noise has a larger relative advantage
    on the v5e at every common configuration. The full grid, including memory limits,
    is in Figure~\ref{fig:f4-full}.}
  \label{fig:f4}
\end{figure*}

On the simplified block, low-rank noise saves more relative to dense noise on the TPU
than on the GPU (Figure~\ref{fig:f4}). Across the same seven feasible configurations,
% generated by experiments/phase2/tb3.py; do not edit
rank-1 generation time is 0.22 times the dense time on the A100 and 0.08 times on the v5e (geometric means).

Regenerating full-rank noise instead trades speed for memory: the seed variant is slower
than dense at every comparable configuration, but fits the complete cost grid.
The full grid and precision comparisons are in Appendix~\ref{sec:systems-ablations}.
The platform difference persists at matched precision: the single-device scaling sweeps
use float32 storage and highest precision on both platforms, and give rank-1/dense time
geometric-mean ratios of 0.43 on A100 and 0.17 on v5e over seven matched
configuration/placement pairs.
These ratios describe the tested workload and implementation; they do not isolate the
hardware mechanism responsible for the difference.

As an implementation check, the library's single-device rank-1 throughput is within 3\% of
EGGROLL's released code on the A100, and 0.68--0.85 times its throughput on the v5e.
The TPU gap was not isolated. Appendix~\ref{sec:refimpl} separates these matched-device
comparisons from the additional throughput available with eight devices.
